\chapter{Konzept und Umsetzung}
\label{cha:umsetzung}

% Erst das Konzept (Architektur, Entwurfsentscheidungen, Alternativen), dann die
% Implementierung (Technologien, zentrale Komponenten, interessante Details).
% Kein Code-Abdruck ganzer Dateien – nur aussagekräftige Ausschnitte mit Erklärung.

\section{Architektur}
\label{sec:architektur}

Abbildung~\ref{fig:architektur} zeigt die Komponenten des Systems und ihre
Schnittstellen. Die Trennung in \ldots{} und \ldots{} erlaubt \ldots

\begin{figure}[htbp]
  \centering
  \missingfigure{Architekturdiagramm einfügen (abbildungen/architektur.png)}
  \caption{Systemarchitektur}
  \label{fig:architektur}
\end{figure}

\section{Entwurfsentscheidungen}
\label{sec:entwurf}

% Für jede wichtige Entscheidung: Alternativen, Kriterien, Begründung.
Für die Datenhaltung wurden zwei Alternativen betrachtet: \ldots{} Gewählt wurde
\ldots, weil \ldots

\section{Implementierung}
\label{sec:implementierung}

Die Umsetzung erfolgte in \ldots{} Listing~\ref{lst:beispiel} zeigt die zentrale
Funktion zur \ldots

% Variante 1: Code direkt aus einer Datei einbinden (bleibt automatisch aktuell)
\lstinputlisting[language=Python, caption={Zentrale Verarbeitungsfunktion}, label={lst:beispiel}, firstline=1, lastline=20]{code/beispiel.py}

% Variante 2: Code im Text
\begin{lstlisting}[language=Python, caption={Aufruf der Funktion}, label={lst:aufruf}]
ergebnis = verarbeite(daten, schwellwert=0.5)
print(ergebnis)
\end{lstlisting}

\section{Herausforderungen}
\label{sec:herausforderungen}

% Was war schwierig und wie wurde es gelöst? Das zeigt Reflexionsfähigkeit.
Eine besondere Herausforderung stellte \ldots{} dar. \todo{Beispiel ergänzen}
